\section{Go Fuzzing for Cryptographic Primitives}


\subsection{The Go Fuzzing Framework}

Go 1.18 introduced native fuzzing in the \texttt{testing} package. A fuzz test is a
function that accepts \texttt{*testing.F} and calls \texttt{f.Fuzz}:

\begin{lstlisting}[language=Go,caption={Basic Go fuzz test for NTT round-trip}]
func FuzzNTTRoundTrip(f *testing.F) {
    // Seed corpus: known interesting inputs
    f.Add([]byte{0, 0, 0, 0})
    f.Add([]byte{255, 255, 255, 255})

    f.Fuzz(func(t *testing.T, data []byte) {
        if len(data) < PolyDegree*4 {
            t.Skip()
        }

        // Decode input as polynomial coefficients
        poly := decodePoly(data[:PolyDegree*4])

        // Apply NTT then INTT
        nttPoly := NTT(poly)
        recovered := INTT(nttPoly)

        // Property: round-trip is identity
        for i := range poly {
            if recovered[i] != poly[i] {
                t.Fatalf("NTT round-trip failed at index %d: "+
                    "got %d, want %d", i, recovered[i], poly[i])
            }
        }
    })
}
\end{lstlisting}

\subsection{Running Go Fuzz Tests}

\begin{lstlisting}[language=bash,caption={Running Go fuzz tests}]
# Run fuzz test for 60 seconds
go test -fuzz FuzzNTTRoundTrip -fuzztime 60s

# Run with race detector (slower but catches data races)
go test -fuzz FuzzNTTRoundTrip -fuzztime 60s -race

# Run with address sanitizer (catches memory bugs in cgo)
CGO_ENABLED=1 go test -fuzz FuzzNTTRoundTrip -fuzztime 60s \
    -asan

# Reproduce a failing input from the corpus
go test -run FuzzNTTRoundTrip/corpus_entry_name
\end{lstlisting}

\subsection{Coverage-Guided Mutation}

Go's fuzzer uses coverage-guided mutation: it instruments the code to track which
branches are taken, then mutates inputs to explore new branches. This is critical for
cryptographic code where interesting behavior often hides behind conditional branches
(e.g., modular reduction triggered only when a coefficient exceeds $q$).

The fuzzer maintains a \emph{corpus}: a set of inputs that collectively maximize branch
coverage. New inputs are added to the corpus only if they cover a previously uncovered
branch.

\subsection{Properties for Polynomial Arithmetic}

\begin{lstlisting}[language=Go,caption={Property-based fuzz tests for polynomial arithmetic}]
// Commutativity: a * b == b * a
func FuzzPolyMulCommutative(f *testing.F) {
    f.Fuzz(func(t *testing.T, seed1, seed2 []byte) {
        a := polyFromSeed(seed1)
        b := polyFromSeed(seed2)

        ab := PolyMul(a, b)
        ba := PolyMul(b, a)

        requirePolyEqual(t, ab, ba)
    })
}

// Distributivity: a * (b + c) == a*b + a*c
func FuzzPolyMulDistributive(f *testing.F) {
    f.Fuzz(func(t *testing.T, s1, s2, s3 []byte) {
        a := polyFromSeed(s1)
        b := polyFromSeed(s2)
        c := polyFromSeed(s3)

        lhs := PolyMul(a, PolyAdd(b, c))
        rhs := PolyAdd(PolyMul(a, b), PolyMul(a, c))

        requirePolyEqual(t, lhs, rhs)
    })
}

// NTT preserves multiplication:
// NTT(a * b) == NTT(a) . NTT(b) (pointwise)
func FuzzNTTHomomorphism(f *testing.F) {
    f.Fuzz(func(t *testing.T, s1, s2 []byte) {
        a := polyFromSeed(s1)
        b := polyFromSeed(s2)

        // Multiply then transform
        ab := PolyMul(a, b)
        nttAB := NTT(ab)

        // Transform then pointwise multiply
        nttA := NTT(a)
        nttB := NTT(b)
        nttProd := PointwiseMul(nttA, nttB)

        requirePolyEqual(t, nttAB, nttProd)
    })
}
\end{lstlisting}

\subsection{FHE Noise Budget Properties}

FHE (Fully Homomorphic Encryption) ciphertexts carry a \emph{noise budget}: each
operation increases noise, and when the budget is exhausted, decryption fails. We fuzz
test noise properties:

\begin{lstlisting}[language=Go,caption={FHE noise budget fuzz test}]
func FuzzFHENoiseGrowth(f *testing.F) {
    f.Fuzz(func(t *testing.T, plaintext1, plaintext2 uint64) {
        // Constrain to valid plaintext space
        p1 := plaintext1 % PlaintextModulus
        p2 := plaintext2 % PlaintextModulus

        // Encrypt
        ct1 := Encrypt(publicKey, p1)
        ct2 := Encrypt(publicKey, p2)

        // Add ciphertexts
        ctSum := Add(ct1, ct2)

        // Noise must not exceed budget
        noise := EstimateNoise(ctSum, secretKey)
        if noise > NoiseBudget {
            t.Fatalf("noise budget exceeded after addition: "+
                "noise=%d, budget=%d", noise, NoiseBudget)
        }

        // Decrypt must succeed
        result := Decrypt(secretKey, ctSum)
        expected := (p1 + p2) % PlaintextModulus
        if result != expected {
            t.Fatalf("decryption failed: got %d, want %d",
                result, expected)
        }
    })
}
\end{lstlisting}

%% ============================================================
